
\clearpage
\section{Holmes and Yarhi-Milo: The Psychological Logic of Peace Summits}

\subsection{Empathy as a mechanism of diplomatic bargaining}

Marcus Holmes and Keren Yarhi-Milo ask why some high-level peace summits succeed while others fail even when the material and strategic setting is broadly similar. Their answer highlights \emph{empathy}: the ability to understand another person's cognitive and affective state without necessarily sharing it, approving it, or sympathizing with it. Empathy therefore has no necessary moral content. It can support cooperation, but it can also be used strategically to understand an opponent's priorities, intentions, and likely behavior.

The authors distinguish cognitive empathy, which concerns understanding another person's perspective and intentions, from affective empathy, which concerns understanding and responding to another person's emotional state. Empathy is also related to, but distinct from, trust. It is difficult to judge whether someone is trustworthy without understanding how that person sees the situation, yet empathy can also reveal that an interlocutor is untrustworthy.

In negotiation, empathy helps actors distinguish stated positions from underlying interests, identify what the other side can actually concede, and determine whether negotiations are being conducted in good faith. This is essential for discovering a zone of possible agreement. A negotiator who cannot understand the counterpart's perspective may treat a symbolic demand as material, misread a concession, or fail to recognize a trade that would benefit both sides.

\subsection{Conveyed empathy and face-to-face interaction}

A central contribution of the article is the distinction between possessing empathic capacity and \emph{conveying} it. Empathic capacity may be a relatively stable personal trait, but it has little diplomatic effect unless the other side perceives it. Leaders infer empathy from expressive behavior in face-to-face interaction: tone, facial expression, posture, eye contact, gestures, mimicry, emotional reactions, and language. These cues can provide information that formal diplomatic statements do not.

This yields the first hypothesis: leaders use beliefs about another leader's empathic capacity as information about whether that leader intends to negotiate in good or bad faith. All else equal, perceived good-faith intentions make agreement more likely, whereas perceived bad faith makes non-agreement more likely. Emotional expression is therefore not simply ``noise'' to be removed from diplomacy; it can communicate sincerity, priorities, and the importance attached to particular issues.

\subsection{Relational empathy and mediation}

In entrenched conflicts, the protagonists may be unable or unwilling to empathize directly. Holmes and Yarhi-Milo therefore introduce \emph{relational empathy}. A mediator can establish an empathic relationship with each side separately and then transmit each protagonist's perspective to the other. If both parties believe that the mediator understands their interests, fears, and political constraints, they are more willing to reveal genuine priorities. The mediator can then correct misunderstandings and identify a bargaining range that the disputants cannot locate themselves.

This produces the second hypothesis: mediators who understand both sides and credibly convey each side's perspective to the other are more likely to generate relational empathy and reach agreement. This role is broader than simple mediator trustworthiness. A mediator can be imperfectly trusted yet still be regarded as someone who understands both parties and can interpret their positions accurately.

\subsection{Camp David 1978: Carter's relational empathy}

The authors compare the successful 1978 Camp David summit among Jimmy Carter, Anwar Sadat, and Menachem Begin with the failed 2000 summit involving Bill Clinton, Yasser Arafat, and Ehud Barak. Both involved high-stakes territorial disputes, deep mistrust, and a Democratic U.S. president acting as mediator, making the contrast useful for assessing the proposed mechanism.

At Camp David I, Carter initially hoped that direct interaction would improve relations between Sadat and Begin. Instead, their meetings quickly became hostile, and the two leaders could barely negotiate face to face. Carter therefore separated them and became the principal channel between the two sides. His role was not limited to carrying proposals. He repeatedly explained each leader's political constraints, motivations, and private flexibility to the other. He reminded Sadat of Begin's ideological and domestic constraints and explained to Begin the political risks Sadat had accepted by visiting Jerusalem and representing broader Arab interests.

Carter was effective because both leaders increasingly believed that he understood them. With Sadat, this relationship developed relatively easily. With Begin it was harder, but emotionally significant interactions -- including the Gettysburg visit and Carter's personal appeal concerning their grandchildren near the end of the summit -- strengthened Begin's sense that Carter understood the human consequences of war and peace.

This relational understanding helped Carter identify the actual bargaining structure. He concluded that Sadat's central objective was the return of Sinai and that Begin might ultimately accept withdrawal and the removal of settlements there, whereas pressing Begin for comparable concessions on the West Bank, Gaza, or Jerusalem was unlikely to succeed. Empathy thus helped Carter distinguish core interests from bargaining style and public rhetoric and to identify a package that other participants had not initially recognized as feasible.

The authors do not claim that empathy alone caused the agreement. Strategic interests, bargaining power, and domestic politics all mattered. Their point is that these factors had to be interpreted and translated into a workable deal, and Carter's empathic mediation made that translation possible.

\subsection{Camp David 2000: failure to build relational empathy}

Camp David II displayed the opposite pattern. Barak and Arafat entered with different expectations and deep suspicion. Their direct interpersonal relationship was poor, and Barak often avoided direct contact with Arafat. Palestinian negotiators interpreted aspects of Barak's behavior as evidence that he did not understand or respect their political needs, while Israeli and American participants often viewed Arafat as evasive and unwilling to state clearly what he would accept.

The absence of direct empathy could have been compensated for by an effective mediator, but Clinton did not create the same relational bridge that Carter had built. Although Clinton was widely regarded as an empathic politician, empathy is contextual and must be perceived by the specific interlocutor. Palestinian participants increasingly believed that Clinton and his team understood Israeli constraints better than Palestinian ones, particularly on Jerusalem, sovereignty, and the political meaning of concessions. As a result, Clinton was less able to serve as a trusted interpreter of each side's underlying interests.

The summit therefore suffered from a relational and informational problem. The participants struggled to distinguish genuine bottom lines from bargaining positions, and no mediator successfully convinced both parties that the other's interests were being accurately understood. Even if Arafat was never prepared to sign a final agreement, the authors argue that the theory still applies: his failure to convey empathic engagement would have signaled the absence of good-faith bargaining. If he was willing to bargain, the failure of relational empathy still reduced the chance that a mutually acceptable package would be found.

\subsection{Conclusion}

Holmes and Yarhi-Milo argue that diplomatic outcomes cannot be explained only by material power, fixed interests, or formal bargaining strategies. Face-to-face interactions influence how leaders interpret those interests and whether they believe counterparts are sincere. Empathy can reveal priorities, clarify intentions, and help actors recognize agreements that would otherwise remain hidden. Where direct empathy is impossible, a skilled mediator can create relational empathy by understanding both sides and translating their perspectives to one another.

Empathy does not guarantee a fair or durable agreement, and it cannot create a bargaining range when none exists. Its importance is instead causal and informational: it affects whether leaders correctly perceive the bargaining opportunities that do exist. The article therefore integrates individual psychology and social relations into the study of diplomacy and shows why mediator behavior can be decisive even when structural conditions remain unchanged.

\section{Kertzer and Tingley: Political Psychology in International Relations -- Beyond the Paradigms}

\subsection{The transformation of political psychology in IR}

Joshua Kertzer and Dustin Tingley review the changing role of political psychology in international relations (IR). Their organizing claim is that psychological research in IR has changed because the broader discipline has changed. Earlier reviews were often structured around grand theoretical paradigms such as realism, liberalism, and constructivism, asking what psychology could contribute to each. Contemporary IR is less dominated by these ``isms,'' less centered exclusively on interstate war, and more methodologically eclectic. Political psychology has followed the same trajectory: it is increasingly problem driven, spans a wider range of substantive areas, and uses a broader set of methods.

The authors identify several reasons for renewed interest in psychology. First, major political developments generated questions that invited psychological explanations: terrorism and radicalization after 9/11, intelligence and planning failures surrounding Iraq and Afghanistan, public backlash against globalization and trade, and renewed attention to individual leaders. Second, IR increasingly shifted toward micro-level approaches and causal mechanisms. Even macro-level theories require assumptions about individual decision making, and scholars interested in mechanisms therefore often encounter psychological processes. The spread of inexpensive survey experiments also made it easier to investigate public attitudes relevant to foreign policy. Third, the traditional opposition between psychology and rational choice weakened. The rise of behavioral economics demonstrated that psychologically realistic assumptions could be incorporated into models of strategic choice rather than treated simply as irrational deviations from them.

This change means that psychology is no longer best understood as a ``theory of errors.'' Psychological mechanisms can explain systematic preferences, emotions, identities, learning processes, and social interactions, not merely mistakes relative to a rational benchmark. At the same time, much contemporary IR work uses psychological ideas without explicitly labeling itself political psychology.

\subsection{A data-driven map of the field}

Rather than relying only on a subjective literature review, Kertzer and Tingley analyze author-selected classifications for manuscripts submitted to \emph{International Studies Quarterly} between 2013 and early 2017. Their figures show both the methodological diversity of political-psychological research and its uneven distribution across substantive topics.

Methodologically, political psychology submissions are much more likely than other IR submissions to use experiments and surveys. Experiments are particularly distinctive: roughly 30\% of political psychology submissions used experimental methods, compared with only a small fraction of submissions overall. Yet political psychology is not methodologically narrow. Case studies remain common, and researchers also use statistics, interviews, archival evidence, content analysis, and interpretive techniques. The field therefore combines the laboratory and survey traditions of psychology with the historical and qualitative traditions of IR.

Substantively, political psychology remains strongly represented in international security and foreign policy, reflecting its historical roots in deterrence, perception, and foreign-policy decision making. It is also active in IR theory, political sociology, human rights, identity, contentious politics, diplomacy, and public opinion. The most striking area of overrepresentation is public opinion: a large share of political-psychological submissions study citizens' attitudes toward international issues.

By contrast, psychology remains relatively underrepresented in international political economy (IPE), international organizations, international institutions, international law, foreign investment, development, monetary politics, and global governance. Kertzer and Tingley do not interpret this pattern as evidence that psychology is irrelevant to these topics. Instead, the gaps identify promising areas for future research because institutions and economic policies ultimately depend on human judgments, identities, and preferences.

\subsection{From cold cognition to hot cognition}

The first major new direction is a shift from ``cold'' cognition toward ``hot'' cognition. Earlier psychological research in IR was deeply shaped by the cognitive revolution and the heuristics-and-biases tradition. It emphasized how decision makers process information, simplify complex environments, misperceive probabilities, and rely on frames or reference points. This work remains important, but newer research pays much greater attention to emotion and affect.

Researchers increasingly distinguish between specific emotions, such as fear, anger, anxiety, humiliation, and pride, because these emotions have different effects on risk taking, threat perception, bargaining, ethnic conflict, and political mobilization. This shift partly overlaps with developments in constructivist IR, where scholars have also turned toward emotion, habit, identity, and non-reflective action. The two literatures sometimes study nearly identical phenomena through different intellectual traditions. Kertzer and Tingley therefore encourage more exchange between political psychologists and constructivists, while noting that constructivist work is often more comfortable with collective emotions and social ontologies whereas psychological research is more often grounded in individual-level mechanisms and experimental evidence.

\subsection{A psychologically richer study of public opinion}

The second development is the rise of psychologically informed research on public opinion in foreign policy. Earlier IR work on public opinion often treated foreign-policy attitudes as a relatively separate domain and used specifically foreign-policy orientations such as militarism or internationalism to explain particular preferences. Newer work increasingly connects foreign-policy attitudes to broader psychological traits and value systems.

Studies show, for example, that personal values, moral foundations, punitiveness, ethnocentrism, social identity, and general orientations toward welfare and redistribution can predict preferences about military force, foreign aid, trade, immigration, and international cooperation. This work implies that foreign-policy preferences are not created only when citizens encounter an international issue. They are partly rooted in more general belief systems that organize political and even nonpolitical choices. In this sense, IR public opinion research has become more integrated with the broader study of political behavior.

\subsection{Bridging the study of leaders, elites, and masses}

A third theme concerns the study of leaders. Classical political psychology often used idiographic techniques such as psychobiography, operational codes, cognitive maps, and personality profiles to analyze individual leaders. Recent leader-centered IR scholarship is more often nomothetic: it uses large datasets to study whether leader characteristics, backgrounds, military experience, or institutional positions systematically affect foreign policy.

The two traditions have not always communicated well. One emphasizes dispositions and individual psychology; the other often emphasizes situations, institutions, and experiences. Kertzer and Tingley argue that future research should study their interaction. Leaders' traits can matter differently depending on institutional constraints, political stakes, and strategic environments.

A similar gap separates elite and mass political psychology. Scholars often use different theoretical frameworks for leaders and ordinary citizens, partly because elites are difficult to recruit for experiments. Yet it is an empirical question, not an assumption, whether elites possess fundamentally different cognitive architectures. The growing use of elite and quasi-elite experiments is therefore valuable. Early findings suggest that elites differ from mass publics on some characteristics, such as strategic expertise or feelings of power, but may still respond to many experimental treatments in similar ways. A more unified psychology of elites and masses would allow scholars to identify when political experience changes basic psychological processes and when it does not.

\subsection{Psychology in international political economy}

The authors identify IPE as a major opportunity for expansion. Psychological approaches are surprisingly rare given the growth of behavioral economics. Existing studies show several productive pathways. Social identity and ethnocentrism help explain trade, immigration, and outsourcing preferences. Loss aversion and inequity aversion help explain economic policy choices. Heuristics and bounded rationality can affect investment decisions and judgments by market actors.

Moreover, classic IPE work on ideas and beliefs often rests implicitly on psychological microfoundations. If ideas constrain which options leaders consider, or if experimentation changes leaders' beliefs, then concepts such as mental models, schema, problem representation, and information processing are directly relevant. The authors therefore call for closer integration between IPE and psychology rather than treating economic preferences as automatically derived from objective material interests.

\subsection{The first image reversed}

One of the article's most original labels is the \emph{first image reversed}. Waltz's first image traditionally asks how individuals affect international politics. Kertzer and Tingley highlight a growing literature that reverses the arrow and asks how international politics affects individuals.

Research on combatants and populations exposed to violence finds that war can reshape political participation, social cooperation, attitudes toward negotiation, risk preferences, and preferences for punishment or retribution. These effects can persist long after the original event and may even be transmitted across generations. The first image reversed therefore shifts attention from the causes of international conflict to its psychological consequences.

This research is important for two reasons. For IR, it connects macro-level events such as war and political violence to everyday political behavior. For political psychology, it emphasizes macro-to-micro causation and demonstrates that dispositions themselves can have situational origins. Rather than treating preferences and psychological traits as fixed inputs, scholars can study how international environments create or activate them.

\subsection{Genetic, biological, and evolutionary approaches}

The sixth major development is the rise of genetic, biological, neuroscientific, and evolutionary approaches. Some research examines genetic contributions to political attitudes; other work investigates hormones, neurobiology, facial characteristics, or evolved behavioral tendencies. Evolutionary perspectives have been applied to questions such as coalition behavior, hierarchy, territoriality, status competition, and threat perception.

Kertzer and Tingley stress that this literature should not be interpreted as a return to biological reductionism. Much contemporary work explicitly studies interactions between biological dispositions and social environments. Genetic effects may depend on context; neurological responses can be shaped by experience; and evolutionary arguments generally concern probabilistic tendencies rather than deterministic behavior. These approaches also create opportunities for IR to export ideas to psychology and related sciences rather than merely import findings from them.

\subsection{Conclusion and research agenda}

The review concludes that political psychology in IR is now broader, more methodologically diverse, and less defined by a simple opposition between rationality and irrationality. Emotions, public opinion, leader characteristics, the psychological effects of conflict, and biological approaches have all expanded the field, while IPE and institutional research remain relatively open frontiers.

The authors add two cautions. First, IR scholars should import psychological findings critically. Psychology is itself diverse and contested, and some concepts that political scientists continue to treat as settled -- including classic formulations of groupthink or the fundamental attribution error -- have been revised or challenged within psychology. Selective borrowing can create the false impression of disciplinary consensus.

Second, political psychologists in IR should aim not only to import theories but also to export findings. International politics provides environments involving high stakes, strategic interaction, violence, institutions, and cross-cultural variation that can generate insights valuable to psychology, behavioral economics, and neuroscience. The future of the field therefore lies in genuine interdisciplinary exchange. Political psychology is most useful when integrated with the strategic, institutional, and social contexts in which international behavior occurs, rather than treated as an isolated alternative paradigm.

\clearpage
\section{Levy: Foreign Policy Decision-Making -- The Psychological Dimension}

\subsection{Political psychology and foreign-policy analysis}

Jack S. Levy surveys the psychological study of foreign-policy decision making and argues that political psychology has become increasingly important in IR after a long period in which individual-level explanations occupied an uncertain position. The chapter combines a history of the field with a review of beliefs, information processing, decision models, group dynamics, and crisis behavior. Levy's central methodological message is that psychological explanations should not be treated as complete theories of international politics by themselves. Their value lies in identifying specific mechanisms inside broader political, institutional, and strategic processes.

For much of the postwar period, mainstream IR privileged systemic and state-level explanations. Structural realism deliberately removed individual psychology from its core theory, while liberal theories concentrated on institutions, regimes, and domestic interests. Constructivism created greater potential for incorporating psychology through its attention to identities, meanings, and socialization, but initially did relatively little to develop individual-level psychological mechanisms. Foreign-policy analysis offered a more natural home for psychological theories because it asked how leaders define situations, process information, and choose policies, yet even early decision-making frameworks often treated leaders' worldviews as given rather than explaining their origins.

The decline of rigid paradigm debates and the rise of middle-range, problem-driven research created more space for psychology. Contemporary work extends beyond security policy into public opinion, trade, finance, human rights, international organizations, and other subjects. Methodologically, the newer behavioral wave is distinguished by quantitative datasets, survey experiments, and, increasingly, elite experiments. At the same time, Levy emphasizes an \emph{aggregation problem}: IR outcomes are produced by collective bodies such as governments, ministries, cabinets, and advisory groups, so researchers must explain how individual beliefs and preferences become collective decisions. A related \emph{strategic interaction problem} arises because foreign-policy decisions interact with the choices of other states. Psychology must therefore move from behavioral decision theory toward something closer to behavioral game theory.

\subsection{Beliefs, images, and operational codes}

Beliefs are central because leaders interpret new information through prior conceptions of how the world works. Political actors hold descriptive beliefs about other states and the international environment, causal beliefs about the consequences of alternative policies, normative beliefs about desirable outcomes, and beliefs about themselves. These beliefs help explain why different leaders facing similar structural conditions may choose different policies.

Early work on national images emphasized perceptions of adversary hostility and strength, as well as corresponding self-images. Enemy images often interact with idealized self-images to create mirror-image processes in which each side interprets its own behavior as defensive and the opponent's similar behavior as aggressive. Holsti's ``inherent bad faith'' model illustrates the resulting rigidity: if an adversary is assumed to be fundamentally hostile, its hostile acts confirm its nature while conciliatory acts are reinterpreted as tactical responses to pressure. Such a belief system becomes difficult to falsify and can obstruct conflict resolution.

George's development of the \emph{operational code} provides a broader framework for organizing leader beliefs. Operational codes include philosophical beliefs about the nature of politics, conflict, predictability, and human agency, as well as instrumental beliefs about which strategies, tactics, timing, and levels of risk are effective. These beliefs are often hierarchical, mutually reinforcing, and resistant to change. Extensions such as the ``crisis bargaining code'' focus specifically on how leaders perceive adversaries, escalation processes, and the appropriate sequencing of coercive and accommodative tactics.

Levy also reviews research linking beliefs to leaders' backgrounds and experiences. New datasets on political leaders make it possible to test systematically whether military experience, combat experience, age, childhood conditions, gender, education, and earlier political careers affect conflict behavior. One notable finding in this literature is that military service without combat experience can be associated with greater conflict propensity, while direct combat experience can have different effects depending on regime context. Prior experience also shapes how leaders are perceived by foreign adversaries. Early actions in office can become anchors for reputations of resolve that are subsequently difficult to change.

Expertise matters as well. Experienced leaders may draw more heavily on specialized knowledge and historical analogies, while less experienced leaders may rely more on personal predispositions or trusted advisers. Yet experienced advisory teams cannot automatically substitute for an experienced leader, because leaders must know enough to monitor delegated advice, recognize disagreements, and evaluate the quality of information they receive.

\subsection{The rational benchmark and bounded rationality}

Psychological approaches to decision making have often been framed as deviations from an idealized rational model. Levy summarizes the rational benchmark as value maximization under constraints. Rational preferences should be complete and transitive, actors should be able to make trade-offs among competing goals, and they should consider intertemporal trade-offs between present and future outcomes. They should search for relevant information, generate alternative strategies, estimate the likely consequences of each, assess probabilities independently of desired outcomes, and update beliefs in a Bayesian manner when new evidence appears.

This standard is difficult to satisfy in real political environments. Simon's concept of \emph{bounded rationality} begins from the fact that decision makers operate with limited time, attention, information, and computational capacity. Rather than optimizing over all possible alternatives, they often satisfice, rely on heuristics, simplify complex problems, and use standard procedures. These strategies are not necessarily foolish: they are often adaptive responses to complexity. But they can also generate systematic distortions.

\subsection{Cognitive biases, selective attention, and belief perseverance}

Levy reviews a range of cognitive biases that influence foreign-policy judgment. Leaders often see what they expect to see because prior beliefs guide attention and interpretation. Anchoring makes initial beliefs resistant to later revision. Confirmation bias encourages the search for supportive evidence and the discounting of contradictory evidence. Once a policy has been chosen, decision makers may engage in selective attention that protects the decision from reconsideration. Belief systems often operate more like theories that organize evidence than like passive collections of observations.

The impact of prior beliefs is especially powerful when evidence is ambiguous. Strong beliefs can cause individuals to interpret the same information in opposite ways. These processes help explain intelligence failures and misperceptions of threat, but Levy cautions that apparent cognitive bias may sometimes have structural or political explanations. Analysts should not label behavior psychologically biased without considering alternative mechanisms.

The \emph{fundamental attribution error} is particularly relevant to international politics. Observers tend to attribute an adversary's hostile behavior to its disposition or intentions while explaining their own comparable behavior as a response to circumstances. This actor-observer asymmetry can reinforce security dilemmas: each side interprets its own armament as defensive but the other's as evidence of aggressive character. Leaders may also underestimate the degree to which adversaries are constrained by domestic politics, bureaucratic pressures, or strategic circumstances.

\subsection{Motivated reasoning}

Motivated reasoning differs from ordinary cognitive bias because preferences and emotional needs directly influence how information is processed. People seek consistency between beliefs and desired outcomes, discount evidence that threatens important commitments, and interpret ambiguous information in ways that protect identities, policies, or interests. Cognitive dissonance theory helps explain why leaders may resist information showing that a favored policy has failed.

Motivated reasoning creates serious methodological problems for historical analysis. A leader's inference may reflect psychological protection of a belief, but it might also be a strategic attempt to justify a policy, satisfy domestic constituencies, or obtain bureaucratic support. Analysts therefore need evidence about process rather than inferring motivation from the outcome alone. The debate over intelligence on Iraq illustrates how difficult it can be to separate biased interpretation from political pressure, ambiguous evidence, and genuine uncertainty.

\subsection{Hawkish biases, overconfidence, and learning from history}

Kahneman and Renshon's concept of ``hawkish biases'' groups together several psychological tendencies that systematically favor coercive or conflictual policies. Positive illusions can produce excessive confidence in one's military capabilities and bargaining leverage. The illusion of control encourages leaders to believe that they can manage escalation better than they actually can. Competition neglect leads actors to focus on their own capabilities while insufficiently considering how rivals will adapt. The illusion of transparency encourages the belief that one's intentions are obvious to others, even when they are not.

These biases can interact. Overconfidence may cause leaders to underestimate risks; belief in control can make escalation appear manageable; and failure to model an adversary's response can create a self-reinforcing spiral. Levy also notes that some apparent biases may be endogenous to conflict: militarized confrontation itself can increase confidence, group cohesion, and hardline beliefs.

Historical analogies constitute another major source of foreign-policy beliefs. Leaders frequently interpret new crises by comparing them to salient past events such as Munich, Vietnam, or the Persian Gulf War. Analogies simplify complex situations and offer ready-made causal lessons, but they can mislead when decision makers focus on superficial similarities and neglect differences in context. The availability heuristic helps explain why dramatic, recent, or personally experienced historical episodes are especially influential. Yet researchers face a difficult inference problem: leaders often invoke analogies rhetorically after reaching a policy preference, so it can be hard to determine whether a historical lesson genuinely caused the decision.

\subsection{Psychology and bargaining models of war}

Levy then connects psychology to rationalist bargaining theories of conflict. Fearon's bargaining model argues that war is puzzling because costly fighting should leave a range of peaceful bargains that rational unitary states prefer. Rationalist explanations therefore focus on private information and incentives to misrepresent, commitment problems, and issue indivisibilities.

Psychological research can complement this framework in several ways. Leaders may disagree about relative power or resolve not only because of private information but also because of stereotypes, overconfidence, biased information processing, or culturally shaped assessments. These mechanisms can distort estimates of the bargaining range even when material information is available. Psychological factors can also shape beliefs about credibility and audience costs. A threat intended as a costly signal may not be interpreted by the target in the way a formal model assumes.

Levy's broader point is that formal rationalist theories identify important strategic structures, while psychology can illuminate where beliefs and subjective probabilities inside those structures come from. This is not a reason to abandon bargaining theory but to build richer models of strategic interaction. Psychological mechanisms matter most when incorporated at clearly specified points in the bargaining process.

\subsection{Prospect theory}

Prospect theory is the leading psychological alternative to expected-utility theory in decisions under risk. It has three central elements. First, outcomes are evaluated as gains or losses relative to a \emph{reference point}, not simply in terms of final asset levels. Second, losses loom larger than equivalent gains: \emph{loss aversion} makes actors particularly resistant to giving up what they regard as already theirs. Third, risk orientation varies across domains. Individuals tend to be risk averse in the domain of gains but more willing to accept risk in the domain of losses. Prospect theory also includes nonlinear probability weighting, whereby people tend to overweight small probabilities and underweight moderate or high probabilities.

These features generate several implications for foreign policy. Reference points can be shaped by the status quo, expectations, recent gains or losses, aspirations, historical analogies, or political framing. Because reference points are partly perceptual, leaders may disagree over whether the same objective outcome constitutes a gain or loss. Once an actor perceives itself to be in a loss domain, it may accept substantial risks to recover the reference point. This logic has been used to explain risky interventions, escalation, and resistance to concessions.

Loss aversion can create a status quo bias and an endowment effect: territory, prestige, or political positions become more valuable once possessed. This can make concessions more difficult than equivalent acquisitions are attractive. It can also reinforce deterrence because actors may fight harder to prevent losses than to obtain new gains. In bargaining, two sides may frame the same settlement differently, with each viewing concessions as losses.

Levy stresses, however, that prospect theory is difficult to test rigorously in international politics. Researchers must identify the relevant reference point independently of the outcome they are trying to explain. They must distinguish loss-domain risk taking from risk taking generated by ideology, personality, strategic incentives, or ordinary probability assessments. Probability weighting further complicates the theory because responses differ between known risks and ambiguous uncertainty. Prospect theory is therefore powerful but should not become a catch-all explanation for any risky foreign policy.

\subsection{Time horizons and intertemporal choice}

Foreign-policy decisions routinely involve trade-offs between immediate and future costs and benefits. Standard rationalist models usually represent these choices through a discount factor, assuming exponential discounting. Behavioral research, however, finds that people often exhibit \emph{hyperbolic discounting}: the relative value of a payoff changes more sharply when it moves from the near future to the present than when the same time interval occurs farther in the future. This produces present bias and time-inconsistent preferences.

Political actors may therefore postpone costly reforms, prefer immediate political rewards over larger future benefits, or repeatedly delay actions that they continue to regard as desirable in the abstract. Time horizons also vary across leaders and institutions. Electoral cycles, regime type, age, career incentives, and perceptions of political survival can affect the weight assigned to future outcomes. Levy notes that IR needs better theories of where these time horizons come from rather than simply inserting different discount factors into models.

\emph{Temporal construal theory} adds that people do not merely discount future outcomes more heavily; they also think about distant and near outcomes differently. Near-term events are represented concretely and in terms of feasibility, while distant events are represented more abstractly and in terms of desirability. This can generate excessive optimism about long-term policies because actors focus on attractive objectives while ignoring implementation details. The theory helps explain why states often plan poorly for the later stages of military interventions or occupations. Leaders with long time horizons may be committed to ambitious future goals while underestimating the practical difficulties of achieving them.

\subsection{Groupthink and group decision making}

Janis's theory of \emph{groupthink} is one of the most influential psychological models of foreign-policy groups. Groupthink refers to a concurrence-seeking tendency in cohesive groups, especially under stress, isolation, and high stakes. Members value consensus and group solidarity, suppress doubts, discount contradictory information, and exert pressure on dissenters. Classic symptoms include illusions of invulnerability, collective rationalization, belief in the group's moral superiority, stereotyping of opponents, self-censorship, and the emergence of ``mindguards'' who shield the group from unwelcome information.

Janis used cases such as the Bay of Pigs and Vietnam to argue that groupthink can produce defective decision making, while the Cuban Missile Crisis was often presented as a case in which organizational procedures helped avoid it. Levy reviews substantial criticism of the model. Empirical studies do not consistently show that cohesion itself causes poor decisions. Some research points instead to group polarization, leadership style, institutional incentives, or social conformity as more precise mechanisms. Cohesive groups can also perform well when they possess norms that encourage critical debate.

The continued relevance of groupthink therefore lies less in treating it as a universal syndrome and more in studying specific group processes. Foreign-policy groups vary in hierarchy, diversity, information flow, leadership intervention, and willingness to tolerate dissent. Scholars should identify which features actually produce conformity and poor information search rather than attributing any failed group decision to groupthink after the fact.

\subsection{Crisis decision making}

International crises intensify many psychological problems. They combine high stakes, uncertainty, severe time pressure, information overload, and often fear of war. As stress rises, decision makers can lose cognitive flexibility, narrow their attention, rely more heavily on prior beliefs and stereotypes, and consider fewer alternatives. Moderate stress may sometimes improve focus, but extreme stress increases the likelihood of rigid and simplified thinking.

Crises also affect organizational processes. Decision making often shifts upward toward top leaders, normal bureaucratic procedures may be bypassed, and ad hoc groups may form. Short-term military considerations can dominate longer-term political objectives. Under time pressure, leaders may favor options that preserve perceived control or avoid immediate losses even when those choices create future risks.

Levy emphasizes that crisis behavior cannot be explained by a single psychological mechanism. The same stressful environment can interact differently with leader personality, prior beliefs, institutional arrangements, group structure, and the strategic behavior of the adversary. The analytical task is to specify these interactions and identify when psychological effects are likely to be strongest.

\subsection{Conclusion: integrating psychology into theories of foreign policy}

Levy concludes that the political psychology of foreign policy has advanced enormously. Its contemporary strength lies in combining historically informed theories of beliefs and perception with experimental and quantitative methods. Major areas of progress include emotions, threat perception, signaling, resolve, public opinion, leader characteristics, and behavioral decision theory.

Important limitations remain. External validity is a persistent concern because many experiments rely on student or mass-public samples rather than political elites making high-stakes decisions. Elite experiments are promising, but they cannot easily reproduce real crises. More research is also needed outside security studies, especially on the psychology of economic decision making and international political economy.

Above all, psychological factors must be embedded within wider theories of decision processes. Analysts need to specify where in the causal chain beliefs, biases, emotions, time horizons, or group dynamics influence policy; how individual-level tendencies are aggregated into government choices; and how those choices interact strategically with other actors. Political psychology is therefore neither a replacement for structural or rationalist IR theory nor a residual explanation for mistakes. It is a set of middle-range mechanisms that can make theories of foreign policy more descriptively accurate and causally complete.

\clearpage
\clearpage
\section{Yarhi-Milo: In the Eye of the Beholder -- How Leaders and Intelligence Communities Assess the Intentions of Adversaries}

\subsection{The problem of inferring political intentions}

Keren Yarhi-Milo examines a central problem in international politics: how policymakers infer the long-term political intentions of adversaries. Assessments of whether another state is expansionist, opportunistic, or satisfied with the status quo can determine deterrence, reassurance, armament, alliance, and diplomatic policies. Yet intentions are not directly observable. Decision makers must infer them from actions, capabilities, statements, and personal interactions, all of which can be ambiguous or strategically manipulated.

The article compares two prominent rationalist approaches with Yarhi-Milo's psychological-organizational alternative. The \emph{behavior thesis} argues that observers should infer intentions from costly non-capability actions, such as joining or withdrawing from binding institutions, undertaking or ending military interventions, or signing arms-control agreements. Because these actions impose costs or constrain future choices, they should be more credible than cheap talk. The \emph{capabilities thesis} argues that military capabilities and changes in them provide the most reliable information. Military buildups should signal a greater capacity and possibly intention to behave aggressively, while reductions can reassure.

Yarhi-Milo develops the \emph{selective attention thesis}. The core argument is that information does not enter decision making neutrally. Individual expectations, psychological biases, vividness, and organizational routines affect which indicators observers notice, treat as credible, and use to infer intentions. Consequently, civilian decision makers and intelligence organizations can examine the same adversary and reach different conclusions because they systematically privilege different forms of evidence.

The article defines political intentions in relation to the status quo. An expansionist state actively seeks to increase power and influence beyond existing boundaries. An opportunistic state would like favorable change but pursues it only when opportunities arise and costs appear low. A status quo state seeks primarily to preserve its relative position. The dependent variable is not whether these assessments were objectively correct, but how observers formed them.

\subsection{Subjective credibility: how political leaders select signals}

The first component of selective attention is the \emph{subjective credibility hypothesis}. It challenges the assumption that policymakers automatically recognize costly signals and update rationally from them. Whether a leader treats evidence as credible depends on prior beliefs, causal theories, and the vividness of the information.

Prior expectations are especially important. Hawks who already believe an adversary is hostile are likely to notice and accept evidence of malign intent while discounting reassuring actions. A costly conciliatory action can be interpreted as deception, economic necessity, or a temporary tactical retreat rather than as evidence of changed goals. Doves are more likely to assimilate the same action into a more benign view. This differs from simple Bayesian updating because the psychological argument allows some information to be ignored or reinterpreted so strongly that meaningful updating does not occur.

Leaders also differ in their causal theories. A policymaker who believes domestic unrest produces diversionary aggression may see instability inside an adversary as an indicator of expansionist intentions, whereas another leader who rejects that theory will not. Thus the same objective information can have different significance because observers disagree about the relationship between observable behavior and underlying type.

Finally, decision makers are disproportionately influenced by \emph{vivid} information: evidence that is concrete, personalized, emotionally engaging, and close in time or space. Face-to-face impressions of foreign leaders can therefore outweigh abstract intelligence about military inventories or doctrine. Personal interaction can generate a powerful ``eyewitness'' sense of knowing an adversary's sincerity. The danger is that vividness and reliability are not the same thing. A personal impression may feel highly diagnostic while being easier to manipulate or more vulnerable to bias than systematic intelligence data.

\subsection{Organizational expertise and intelligence assessments}

The second component of selective attention is the \emph{organizational expertise hypothesis}. Intelligence organizations face a different information environment from political leaders. They are bureaucracies tasked with collecting, measuring, and comparing data. Military capabilities are especially attractive indicators because they are relatively observable, quantifiable, and trackable over time. Weapons, force deployments, and military expenditures can be counted and incorporated into standardized assessments.

Over time, this repeated monitoring creates organizational expertise and also a form of narrowness. Intelligence agencies become highly skilled at analyzing what they routinely collect, and they are therefore inclined to use that information when asked to infer political intentions. The claim is not that intelligence analysts believe capabilities mechanically determine intentions. Rather, in an inherently difficult task such as estimating long-term political goals, organizations naturally rely on the indicators for which they possess the deepest institutional knowledge.

This differs from the capabilities thesis in causal logic. The capabilities thesis says observers should use military power because it objectively reveals or constrains intentions. The organizational expertise thesis says intelligence communities use military power because their routines, resources, and accumulated expertise make it the most available and institutionally legitimate evidence.

\subsection{Capabilities, behavior, and observable predictions}

The capabilities thesis links intentions to changes in the adversary's military power. A large military buildup, especially by an already powerful adversary, should be interpreted as a costly signal of more hostile intentions. A freeze or decline in capabilities should reassure. If this thesis is correct, assessments should change when perceived capabilities change, and observers should explicitly cite those changes as the reason for revising their beliefs.

The behavior thesis focuses instead on costly actions that are not simply changes in the quantity of military power. Joining a constraining institution, accepting arms-control obligations, withdrawing from a foreign intervention, or otherwise paying a meaningful cost should reveal benign intent; intervention, institutional withdrawal, treaty violation, or other costly revisionist behavior should reveal hostility. The timing of belief change is important. If signing an arms-control agreement itself changes perceptions before capabilities are actually reduced, that supports the behavior thesis. If perceptions change only after material reductions occur, that supports the capabilities thesis.

Yarhi-Milo summarizes these competing predictions in a table and tests them through three historical cases: the collapse of U.S.-Soviet detente under Jimmy Carter, the end of the Cold War during Ronald Reagan's second term, and British assessments of Nazi Germany before World War II. The research draws on more than 30,000 archival documents, intelligence reports, and interviews. The design compares political leaders with intelligence communities and uses process tracing to identify not only whether perceptions changed but which indicators actors explicitly used to justify the change.

\subsection{The collapse of detente, 1977--1980}

The Carter administration provides a strong test because Soviet behavior generated multiple potentially informative signals. Soviet military capabilities were growing and modernizing. Moscow intervened in a number of Third World conflicts, culminating in the invasion of Afghanistan. At the same time, the Soviet Union signed SALT II, an action that could be interpreted as a costly institutional and arms-control signal.

Carter's senior decision makers began with different priors. National Security Adviser Zbigniew Brzezinski was more suspicious of Soviet intentions than Carter or Secretary of State Cyrus Vance. Early in the administration, Carter and Vance generally viewed Soviet intentions as at most opportunistic, while Brzezinski was more alert to the possibility of a coordinated expansionist strategy.

As Soviet involvement in Africa increased, Brzezinski gradually interpreted the pattern as evidence of expansionism. Vance saw much of the same activity as opportunistic rather than part of a grand design. Carter initially resisted Brzezinski's interpretation and continued to seek cooperation. These differences are difficult to explain through the objective costliness of Soviet actions alone, because the officials disagreed about what the same behavior meant.

The invasion of Afghanistan became a decisive moment for Carter. His beliefs shifted sharply toward the view that Soviet intentions were expansionist. Yarhi-Milo argues that the change was not driven mainly by a new calculation of the military balance. The invasion was vivid, dramatic, and personally meaningful. Carter interpreted it as revealing a qualitatively different Soviet willingness to use force and as invalidating his prior expectations about detente. His personal reaction to Soviet leader Leonid Brezhnev also mattered: Carter had previously believed that personal interaction had helped him understand Brezhnev, which made the invasion feel like evidence that this trust had been misplaced.

Vance responded differently. He opposed the invasion but did not infer from it an overarching Soviet plan for global expansion. His interpretation remained closer to the view that Soviet actions were situational and opportunistic. Brzezinski, by contrast, saw Afghanistan as confirming beliefs he had increasingly held for some time. The variation among officials therefore strongly supports the subjective credibility mechanism: the same costly events were selected and interpreted through different priors and theories.

The behavior thesis receives partial support because interventions clearly mattered, but their impact was not uniform or automatic. The capabilities thesis is weaker at the political-leader level. Although decision makers knew that Soviet military capabilities were improving, they rarely used the military buildup itself as the primary basis for judging Soviet political intentions.

\subsection{The U.S. intelligence community during the Carter period}

The intelligence community displayed a different pattern. National Intelligence Estimates concentrated heavily on Soviet military capabilities and the strategic balance. Analysts repeatedly linked growing Soviet power to the possibility of more assertive foreign behavior. Agencies disagreed at times over the precise military balance, but those disagreements were themselves reflected in different judgments about Soviet intentions: analysts who viewed Soviet capabilities as relatively weaker tended to view Soviet objectives as less aggressive, while those who saw a more favorable balance for Moscow inferred greater assertiveness.

This pattern is consistent with both the capabilities thesis and Yarhi-Milo's organizational expertise hypothesis. The intelligence community had developed extensive routines and expertise for measuring forces, weapons, and strategic trends, so it relied on those indicators when faced with the much harder task of estimating intentions. It did not generally base its long-term assessments on vivid personal information of the kind that influenced political leaders.

\subsection{The end of the Cold War, 1985--1988}

The Reagan case provides an especially demanding test for the behavior thesis because Mikhail Gorbachev undertook numerous costly actions that should have reassured observers. During Reagan's first term, the president regarded the Soviet Union as deeply expansionist and ideological. In the second term, however, Reagan and Secretary of State George Shultz gradually came to believe that Soviet intentions had changed. Secretary of Defense Caspar Weinberger remained much more skeptical.

Material capabilities did not offer a straightforward reason for Reagan's change. U.S. assessments did not initially show a dramatic Soviet military decline, and some Soviet modernization continued. Gorbachev's reforms and foreign-policy initiatives were therefore more important to the change in perceptions than the balance of power itself.

Some of Gorbachev's actions were highly costly: acceptance of the Intermediate-Range Nuclear Forces Treaty, willingness to accept intrusive verification and asymmetrical reductions, movement toward withdrawal from Afghanistan, domestic political reforms, greater openness, and reduced interventionism. Reagan and Shultz paid attention to these developments, but Yarhi-Milo argues that their influence was filtered through vivid interpersonal experience. Reagan met Gorbachev repeatedly and came to see him as sincere, pragmatic, and genuinely interested in changing Soviet-American relations. Personal impressions increased the credibility Reagan assigned to Gorbachev's actions. He increasingly referred to Gorbachev's trustworthiness and to the evidence he obtained from private conversations.

Shultz experienced a similar process through his relationship with Soviet Foreign Minister Eduard Shevardnadze. Repeated face-to-face contact produced a sense that the Soviet leadership was changing in ways not captured by abstract indicators alone. Weinberger, who had much less personal interaction with Gorbachev and began from more hawkish expectations, discounted many of the same actions. He interpreted reforms and arms-control initiatives as tactical measures designed to strengthen rather than transform Soviet power.

This divergence again supports the subjective credibility hypothesis. Costly signals were not self-interpreting. Their meaning depended on the prior beliefs and personal experiences of the observer. The behavior thesis explains part of the case because costly Soviet actions clearly mattered, but it cannot explain why some officials updated strongly while others did not.

\subsection{The intelligence community and Gorbachev}

The U.S. intelligence community revised its assessments much more slowly than Reagan and Shultz. It continued to view Soviet intentions through the lens of military capabilities and institutionalized estimates. Even as Gorbachev introduced reforms, intelligence assessments often interpreted them as attempts to strengthen the Soviet state and preserve its long-term capacity for competition. Analysts were cautious about treating political reforms or diplomatic initiatives as proof of fundamentally changed goals.

The lag illustrates organizational expertise. The intelligence community had enormous knowledge about military capabilities but less institutional capacity to interpret the political meaning of personal interactions or rapid ideological change. Its procedures encouraged continuity and made dramatic revisions difficult until material and behavioral trends became overwhelming. This did not make intelligence analysis irrational; it reflected the information its organizations were structured to collect and validate.

\subsection{British assessments of Nazi Germany, 1934--1939}

The interwar case provides variation in both military capabilities and costly behavior. Germany rearmed rapidly, withdrew from the League of Nations and the Disarmament Conference, remilitarized the Rhineland, absorbed Austria, demanded the Sudetenland, and finally occupied the remainder of Czechoslovakia. If military capabilities or objectively costly behavior mechanically determined perceived intentions, British leaders should have converged rapidly on a hostile assessment.

They did not. Officials with different prior beliefs drew sharply different conclusions. Robert Vansittart, who began with a hawkish interpretation of Nazi Germany, viewed German rhetoric and behavior as evidence of a broad revisionist and expansionist project. Other officials, including Neville Chamberlain and Ambassador Nevile Henderson, initially interpreted Hitler's goals as more limited and potentially satisfiable. The same German actions could therefore be categorized as proof of expansionism or as understandable efforts to revise an unfair postwar settlement.

Personal interaction was central to Chamberlain's judgments. His meetings with Hitler during the Sudeten crisis gave him a vivid sense that the German leader's territorial demands were limited and that Hitler could be trusted once those demands were met. Chamberlain gave considerable weight to Hitler's personal assurances even though other evidence pointed toward broader revisionism. This is exactly the type of inference predicted by the subjective credibility hypothesis: concrete interpersonal impressions can outweigh abstract or historical indicators.

The German occupation of the rest of Czechoslovakia in March 1939 finally transformed Chamberlain's view. The action mattered partly because it was costly and violated earlier commitments, which offers support for the behavior thesis. Yet its psychological impact was also highly personal: it directly contradicted assurances Chamberlain believed Hitler had given him. The event therefore destroyed not simply a policy expectation but Chamberlain's personal confidence in his ability to read Hitler. He and Foreign Secretary Halifax increasingly concluded that Germany sought domination rather than limited revision.

Even here, the behavior thesis is incomplete because officials did not agree in advance about which German actions were diagnostic. Some regarded withdrawal from international institutions, the Rhineland, or Austria as clear signals of expansionism; others discounted them. Costliness alone did not fix the meaning of the signal.

\subsection{British military intelligence}

British intelligence assessments again differed from those of political leaders. The Chiefs of Staff and military intelligence concentrated heavily on the German military balance and rearmament. Their task was to estimate what Germany could do, when it would be ready for war, and how British forces compared. This emphasis sometimes led to counterintuitive conclusions about intentions. Even as German political behavior became increasingly aggressive, assessments of military readiness could imply that Germany was unlikely to initiate a major war immediately because it was not yet prepared to defeat Britain and its allies.

The case therefore offers support for the organizational expertise hypothesis. Intelligence organizations repeatedly used capabilities because those were the indicators they were structured to measure. Political leaders, meanwhile, relied more on behavior, ideology, diplomatic interaction, and personal impressions. The capabilities thesis alone cannot explain why civilian leaders did not consistently use the dramatic German buildup as their primary indicator of Hitler's political goals.

\subsection{Conclusions for signaling theory and intelligence}

Across the three cases, Yarhi-Milo finds the strongest support for selective attention. Decision makers did not simply collect all available costly signals and update proportionally. Their prior beliefs shaped which evidence they treated as credible; vivid personal information frequently had disproportionate influence; and different officials could draw opposite inferences from the same costly action. Carter's response to Afghanistan, Reagan's changing view of Gorbachev, and Chamberlain's changing view of Hitler all illustrate how personally salient events can trigger major belief revision.

Intelligence communities followed a more stable and capability-centered logic. Their organizational expertise in military inventories, force balances, and observable trends made those indicators central to their assessments. This generated systematic differences between political leaders and professional intelligence analysis. Neither perspective was necessarily always superior. Leaders' interpersonal judgments could reveal important political changes that intelligence organizations missed, but they were also vulnerable to manipulation and bias. Intelligence assessments could be disciplined and data rich, yet overly dependent on what agencies were best equipped to measure.

The findings qualify rationalist costly-signaling theories. Yarhi-Milo does not reject the idea that costly actions can convey information. Some clearly did. The problem is that the theory often leaves the receiver of the signal underdeveloped. A signal has no political effect merely because it is objectively costly; observers must notice it, classify it as relevant, and interpret what it means. Different costly signals can also point in opposite directions, and rationalist accounts often do not specify which one the observer will privilege.

For policy, the implication is that states cannot assume their intended signals will be understood as intended. A military reduction, treaty, intervention, or public assurance may be filtered through the adversary's existing beliefs and organizational routines. Leaders seeking reassurance or deterrence need to understand not only the objective cost of a signal but also the psychology and institutional position of the audience receiving it. Intelligence organizations, meanwhile, should recognize that their expertise can create blind spots when political intentions change faster than military capabilities.

The article therefore reframes the study of intentions from a problem of identifying the objectively best signal to a problem of perception and attention. Intentions are ``in the eye of the beholder'' because evidence acquires meaning through individual expectations and organizational practices. Any adequate theory of signaling and threat assessment must explain both what states do and how specific observers decide what those actions mean.
